\epanalipsi{Β2}

\tmima{Λεξικό του κεφαλαίου}

\renewcommand{\arraystretch}{1.5}
\begin{tabular*}{\linewidth}{@{\extracolsep{\fill}} >{\raggedright\arraybackslash}p{0.28\linewidth} >{\raggedright\arraybackslash}p{0.62\linewidth}}
\toprule
\textbf{Λέξη ή σύμβολο} & \textbf{Με μία πρόταση}\\
\midrule
συμμετρικό ως προς άξονα $\xi$ & Το $\mathrm{Α}'$ με τον $\xi$ μεσοκάθετο του $\mathrm{ΑΑ}'$. Τα συμμετρικά σχήματα είναι ίσα.\\
άξονας συμμετρίας & Ευθεία που χωρίζει το σχήμα σε δύο σχήματα που καθρεφτίζονται.\\
μεσοκάθετος & Η κάθετη στο μέσο ενός τμήματος· κάθε σημείο της ισαπέχει από τα άκρα.\\
συμμετρικό ως προς σημείο $\mathrm{Ο}$ & Το $\mathrm{Α}'$ με το $\mathrm{Ο}$ μέσο του $\mathrm{ΑΑ}'$ (στροφή $180^\circ$).\\
κέντρο συμμετρίας & Σημείο $\mathrm{Ο}$ γύρω από το οποίο το σχήμα, με στροφή $180^\circ$, πέφτει πάνω στον εαυτό του.\\
τέμνουσα & Ευθεία που τέμνει δύο άλλες ευθείες.\\
εντός εναλλάξ & Ίσες γωνίες, όταν οι ευθείες είναι παράλληλες.\\
εντός επί τα αυτά & Παραπληρωματικές γωνίες, όταν οι ευθείες είναι παράλληλες.\\
\bottomrule
\end{tabular*}

\tmima{Ασκήσεις επανάληψης}

\ex Βρες το συμμετρικό του σημείου $\mathrm{Α}$, που απέχει $3{,}4$ cm από τον άξονα $\xi$. Πόσο απέχουν τα $\mathrm{Α}$ και $\mathrm{Α}'$;

\ex Πόσους άξονες συμμετρίας έχουν: ισόπλευρο τρίγωνο, ορθογώνιο μη τετράγωνο, ρόμβος, κύκλος;

\ex Σημείο $\mathrm{Μ}$ της μεσοκαθέτου του $\mathrm{ΑΒ}$ απέχει $5$ cm από το $\mathrm{Α}$. Πόσο απέχει από το $\mathrm{Β}$;

\ex Ένα τμήμα $\mathrm{ΑΒ}=10$ cm. Πόσο απέχει η μεσοκάθετός του από το $\mathrm{Α}$;

\ex Το τρίγωνο $\mathrm{ΑΒΓ}$ έχει $\mathrm{ΑΒ}=\mathrm{ΑΓ}$. Ποιος είναι ένας άξονας συμμετρίας του;

\ex Βρες το συμμετρικό του $\mathrm{Α}(-2)$ στον άξονα των αριθμών ως προς το $0$ και ως προς το $1$.

\ex Σημείο $\mathrm{Α}$ απέχει $7$ cm από το $\mathrm{Ο}$. Πόσο απέχουν τα $\mathrm{Α}$ και $\mathrm{Α}'$ (συμμετρικό ως προς $\mathrm{Ο}$);

\ex Ποια κεφαλαία γράμματα από τα Α, Β, Η, Ν, Σ, Τ, Φ έχουν άξονα και ποια κέντρο συμμετρίας;

\ex Δύο παράλληλες ευθείες τέμνονται από τρίτη και μία γωνία είναι $80^\circ$. Πόσες γωνίες είναι $80^\circ$ και πόσες $100^\circ$;

\ex Οι εντός επί τα αυτά γωνίες δύο παραλλήλων είναι $2x$ και $x+60^\circ$. Βρες το $x$.

\ex Δύο εντός εναλλάξ γωνίες είναι $64^\circ$ και $67^\circ$. Είναι παράλληλες οι ευθείες; Γιατί;

\ex Η μεσοκάθετος και η διαγώνιος: Σε ρόμβο $\mathrm{ΑΒΓΔ}$ η $\mathrm{ΒΔ}$ είναι ποιας διαγωνίου μεσοκάθετος;

\ex Ένα ορθογώνιο $12\times5$ έχει διαγώνιο $13$. Πού βρίσκεται το κέντρο συμμετρίας και πόσο απέχει από κάθε κορυφή;

\ex Πόσο είναι το μήκος του συμμετρικού τμήματος $\mathrm{Α}'\mathrm{Β}'$ ενός τμήματος $\mathrm{ΑΒ}=8{,}5$ cm; Τι σχέση έχει με το $\mathrm{ΑΒ}$ αν η συμμετρία είναι ως προς σημείο εκτός της ευθείας $\mathrm{ΑΒ}$;

\ex \textbf{Ερωτήσεις αυτοαξιολόγησης (συμμετρίες).} Σωστό ή λάθος; α) Τα συμμετρικά σχήματα ως προς άξονα είναι ίσα. \quad β) Η μεσοκάθετος περνά από το μέσο του τμήματος. \quad γ) Το ισοσκελές τρίγωνο έχει κέντρο συμμετρίας. \quad δ) Ο κύκλος έχει άπειρους άξονες.

\ex \textbf{Ερωτήσεις αυτοαξιολόγησης (παράλληλες).} Σωστό ή λάθος; α) Οι εντός εναλλάξ γωνίες δύο παραλλήλων είναι ίσες. \quad β) Οι εντός επί τα αυτά είναι ίσες. \quad γ) Αν οι εντός εναλλάξ είναι ίσες, οι ευθείες είναι παράλληλες. \quad δ) Οι κατακορυφήν γωνίες είναι ίσες.

\tmima{Ο χάρτης της επανάληψης}

Μετά τις λύσεις, δες σε ποιες ασκήσεις δυσκολεύτηκες. Ο πίνακας σου λέει ποια παράγραφο να ξαναδιαβάσεις.

\renewcommand{\arraystretch}{1.5}
\begin{tabular*}{\linewidth}{@{\extracolsep{\fill}} >{\raggedright\arraybackslash}p{0.38\linewidth} >{\raggedright\arraybackslash}p{0.52\linewidth}}
\toprule
\textbf{Αν δυσκολεύτηκες στις} & \textbf{ξαναδιάβασε τις παραγράφους}\\
\midrule
ΕΒ2.1, ΕΒ2.14 & 2.1 (συμμετρία ως προς άξονα)\\
ΕΒ2.2, ΕΒ2.5, ΕΒ2.8 & 2.2 (άξονες συμμετρίας)\\
ΕΒ2.3, ΕΒ2.4, ΕΒ2.12 & 2.3 (μεσοκάθετος)\\
ΕΒ2.6, ΕΒ2.7 & 2.4 (συμμετρία ως προς σημείο)\\
ΕΒ2.13 & 2.5 (κέντρο συμμετρίας)\\
ΕΒ2.9, ΕΒ2.10, ΕΒ2.11 & 2.6 (παράλληλες και τέμνουσα)\\
ΕΒ2.15, ΕΒ2.16 & 2.1 έως 2.6\\
\bottomrule
\end{tabular*}
